\startbuffer[sacc]
    {\emph Saccharomyces cerevisiae}
\stopbuffer

\startbuffer[peni]
    {\emph Penicillium chrysogenum}
\stopbuffer

\startchapter[title={Revisão}]
    A industria de biotecnologia é responsável por produções alternativas de compostos usuais na nossa sociedade, utilizando organismos vivos para produzir alimentos, biocombustiveis \cite[authoryear][ref::Nielsen], farmacos \cite[authoryear][ref::Zhang] e compostos químicos \cite[authoryear][ref::Meadows]. No balanço de sustentabilidade, estas rotas alternativas, utilizando organismos modificados ou não, permitem uma aproximação orgânica e amigavel com o meio ambiente para os processos produtivos envolvidos, promovendo uma transição para uma economia verde, ou seja, reduzindo o saldo de carbono emitido \cite[authoryear][ref::Naseri]. Estas alternativas não são à principio uma degressão em produtividade nas cadeias produtivas mencionadas em prol da sustentabilidade, mas sim, puramente uma evolução em eficiência na alocação de espaço físico, consumo de insumos e geração de subprodutos. A sustentabilidade é mais um ponto positivo que pesa a favor de implementações em grande escala de processos que utilizem técnicas em biotecnologia em suas rotas produtivas.

    \startsection[title={Bioreatores}]
        Bioreatores fornecem ambientes favoráveis para o crescimento de culturas de microrganismos e podem ser otimizados para a produção de determinado produto bioquímico, desde que haja acesso a mecanismos de controle das variáveis físicas envolvidas, o que geralmente é oferecido por reatores projetados especificamente para essa finalidade. A necessidade de um projeto de bioreator decorre da responsabilidade de garantir a reprodutibilidade dos processos bioquímicos desenvolvidos ao longo de uma cadeia produtiva. A possibilidade de controlar algumas variáveis --- ou até mesmo todas aquelas avaliadas como impactantes para a conversão final --- é uma característica indispensável para aplicações industriais de rotas produtivas desenvolvidas em laboratório. A princípio, uma vez determinadas as condições de operação na etapa experimental, basta ajustar as variáveis extensivas do processo para uma escala industrial.

        Segundo \cite[authoryears][ref::Erickson], existem diversas formas de classificar bioreatores, seja pelo tipo de alimentação ou pelos produtos finais obtidos. Existem três tipos principais de alimentação: batelada, alimentação alternada e alimentação contínua. A alimentação em batelada ocorre quando o substrato é introduzido no sistema apenas no início do processo, sem manipulações externas posteriores na quantidade de fonte de carbono. A alimentação alternada consiste em uma estratégia de distribuição da fonte de carbono seguindo um cronograma, diferenciando-se da primeira pela possibilidade de introdução de substrato durante o experimento. Por fim, a alimentação contínua caracteriza-se pela introdução constante de alimento ao longo de todo o processo, representando uma tentativa de reproduzir um sistema estacionário a partir da manutenção da concentração de substrato no bioreator.

        \startplacefigure[title={Ilustração de bioreator genérico}]
            \externalfigure[image/bioreactor.png][height=200pt]
        \stopplacefigure

        Quando tratamos de processos aeróbicos, é necessário fornecer oxigênio para o desenvolvimento da cultura. Nessa perspectiva, boa parte dos reatores destinados ao cultivo de organismos aeróbicos necessita de agitação para uniformizar a concentração de oxigênio em todo o meio, o que é naturalmente dificultado pela baixa solubilidade do oxigênio em água. Isso pode ser realizado por meio de reatores de mistura, colunas de ar dissolvido, membranas ou pelo aumento da área de interface líquido-gás, o que não ocorre necessariamente em processos anaeróbicos ou em sistemas envolvendo culturas mistas \cite[authoryear][ref::Erickson]. Essa última opção tem como proposta estimular a competição entre diferentes variantes presentes com base na seleção natural. Portanto, os tipos de cultura influenciam diretamente no projeto físico do bioreator.

        Outra condição de operação que influencia drasticamente tanto o projeto físico quanto o modelo matemático de um reator é a fase dos reagentes e produtos. Em fase líquida, a cinética bioquímica depende da concentração dos reagentes e, em alguns casos, também da concentração dos produtos. Para muitas reações bioquímicas, essa abordagem cinética pode ser tratada, com boa aproximação, considerando dependência de primeira ordem em baixas concentrações e de ordem zero em concentrações elevadas \cite[authoryear][ref::Michaelis]. Em fase sólida, a cinética bioquímica torna-se mais dependente da disposição espacial do substrato sólido e do contato superficial deste substrato com o microrganismo. Dessa forma, uma solubilidade substancial dos reagentes é essencial para garantir adequada difusão no meio e, consequentemente, a formação do produto de interesse \cite[authoryear][ref::Bhargav]. Em alguns sistemas, utilizam-se bioenzimas para hidrolisar o substrato em reagentes mais solúveis, aumentando a conversão por unidade de tempo.

        O controle das variáveis físicas envolvidas é importante para a predição da resposta do sistema em relação a qualquer perturbação de entrada, bem como para a manipulação do sistema, de forma que os resultados de réplicas experimentais sejam os mais semelhantes possíveis dentro das limitações de sistemas não determinísticos. Entre as variáveis que podem ser medidas em um bioreator estão temperatura, pressão, pH, potência de agitação, fluxo de gases e líquidos, concentração de gases, concentração de biomassa e tensão interfacial \cite[authoryear][ref::Erickson].
    \stopsection

    \startsection[title={Fermentação}]
        A fermentação ocupa posição de destaque na indústria biotecnológica, assumindo protagonismo, ainda hoje, em rotas produtivas de compostos essenciais para a sociedade \cite[authoryear][ref::Metcalfe]. Grande parte das vacinas, medicamentos e biocombustíveis é obtida preferencialmente por processos fermentativos, e sua viabilidade econômica decorre justamente da rota empregada. Caso contrário, métodos alternativos de produção seriam mais trabalhosos, dispendiosos e provavelmente economicamente inviáveis na escala atual.

        Os microrganismos responsáveis pela fermentação podem ser bactérias, leveduras, arqueias, fungos ou células de mamíferos \cite[authoryear][ref::Muller]. Portanto, existe uma ampla variedade de caminhos metabólicos possíveis considerando o conjunto desses organismos. As leveduras, por exemplo, são responsáveis pela produção de etanol, pão, bebidas alcoólicas, penicilina, entre outros produtos. Isso reforça a relevância econômica e social desses microrganismos e de suas aplicações industriais.

        A produção de etanol por fermentação é um dos processos biotecnológicos mais importantes da indústria moderna, sendo amplamente utilizada para a obtenção de biocombustíveis e produtos químicos renováveis \cite[authoryear][ref::Dashko]. Nesse processo, leveduras como \getbuffer[sacc] convertem açúcares provenientes de matérias-primas vegetais, como cana-de-açúcar, milho e biomassa lignocelulósica, em etanol e dióxido de carbono sob condições anaeróbicas \cite[authoryear][ref::Bhargav]. A fermentação alcoólica apresenta grande relevância econômica, especialmente em países como o Brasil, onde o etanol é utilizado em larga escala como combustível automotivo. Além de sua aplicação energética, o etanol também é empregado nas indústrias farmacêutica, alimentícia e química, demonstrando a versatilidade dos processos fermentativos na produção industrial \cite[authoryear][ref::Bhargav].

        A utilização da penicilina possibilitou avanços importantes na medicina moderna, como cirurgias mais seguras, tratamentos hospitalares mais eficazes e aumento da expectativa de vida da população. Atualmente, apesar do surgimento de bactérias resistentes, a penicilina e seus derivados ainda possuem grande importância terapêutica no tratamento de diversas infecções \cite[authoryear][ref::Goldrick]. A obtenção da penicilina ocorre principalmente por meio de processos fermentativos utilizando fungos do gênero \getbuffer[peni]. Durante a fermentação, o fungo produz penicilina como metabólito secundário, que posteriormente é extraído, purificado e processado para uso farmacêutico.

        \startplacefigure[title={Ilustração de levedura}, reference=fig:yeast_ilustration]
            \externalfigure[image/yeast_ilustration.jpg][height=200pt]
        \stopplacefigure

        \startplacefigure[title={Ilustração de fungo}, reference=fig:fungal_ilustration]
            \externalfigure[image/fungal_ilustration.jpg][height=200pt]
        \stopplacefigure

        A \in{Figura}[fig:yeast_ilustration] mostra uma célula de levedura. Nela, é possível observar o núcleo, onde se encontra o material genético principal dos estudos de modificação genética. A mitocôndria também possui material genético, embora a extensão dos estudos relacionados à sua manipulação seja menor quando comparada à do núcleo. Um ponto importante na compreensão das reações químicas de oxirredução \cite[authoryear][ref::Muller] é que elas ocorrem no citoplasma ou são coordenadas na mitocôndria quando associadas a processos aeróbicos.
    \stopsection

    \startsection[title={Fermentação Alcoólica}]
        No que se refere à produção de álcool a partir de fontes de carbono, as leveduras desempenham papel fundamental. Em especial, \getbuffer[sacc] representa um grupo extremamente relevante nesse processo. Do ponto de vista evolutivo, essas leveduras desenvolveram essa capacidade inicialmente em resposta ao surgimento de frutas, fontes naturais de açúcares simples, possibilitando a conversão desses açúcares em álcool tanto em condições aeróbicas quanto anaeróbicas \cite[authoryear][ref::Dashko].

        Em geral, o metabolismo de um microrganismo pode ser modificado por meio da inserção de genes reacionais em estequiometria adequada para maximizar determinado produto ou inibir mecanismos metabólicos específicos. \cite[authoryears][ref::Meadows] trabalhou na modificação do genoma de variantes de \getbuffer[sacc] e demonstrou a possibilidade de aumentar significativamente a seletividade de caminhos metabólicos de interesse, especificamente na produção de precursores de ácidos isoprenoides utilizados no tratamento da malária. Em situações como essa, estudadas por \cite[author][ref::Meadows], o metabolismo torna-se diferente daquele naturalmente encontrado em variantes de \getbuffer[sacc]. Entretanto, ao restringir a análise ao metabolismo nativo genérico, caracterizado principalmente pelo consumo de glicose e pela formação de etanol, é possível visualizar o mecanismo representado na \in{Figura}[fig:alcoholic_fermentation].

        \startplacefigure[title={Mecanismo genérico parcial de fermentação alcoólica}, reference=fig:alcoholic_fermentation]
            \externalfigure[image/alcoholic_fermentation.png][height=200pt]
        \stopplacefigure
    \stopsection
\stopchapter
